% Options for packages loaded elsewhere
\PassOptionsToPackage{unicode}{hyperref}
\PassOptionsToPackage{hyphens}{url}
%
\documentclass[
  11pt,
]{article}
\usepackage{amsmath,amssymb}
\usepackage{iftex}
\ifPDFTeX
  \usepackage[T1]{fontenc}
  \usepackage[utf8]{inputenc}
  \usepackage{textcomp} % provide euro and other symbols
\else % if luatex or xetex
  \usepackage{unicode-math} % this also loads fontspec
  \defaultfontfeatures{Scale=MatchLowercase}
  \defaultfontfeatures[\rmfamily]{Ligatures=TeX,Scale=1}
\fi
\usepackage{lmodern}
\ifPDFTeX\else
  % xetex/luatex font selection
\fi
% Use upquote if available, for straight quotes in verbatim environments
\IfFileExists{upquote.sty}{\usepackage{upquote}}{}
\IfFileExists{microtype.sty}{% use microtype if available
  \usepackage[]{microtype}
  \UseMicrotypeSet[protrusion]{basicmath} % disable protrusion for tt fonts
}{}
\makeatletter
\@ifundefined{KOMAClassName}{% if non-KOMA class
  \IfFileExists{parskip.sty}{%
    \usepackage{parskip}
  }{% else
    \setlength{\parindent}{0pt}
    \setlength{\parskip}{6pt plus 2pt minus 1pt}}
}{% if KOMA class
  \KOMAoptions{parskip=half}}
\makeatother
\usepackage{xcolor}
\usepackage[margin=1in]{geometry}
\setlength{\emergencystretch}{3em} % prevent overfull lines
\providecommand{\tightlist}{%
  \setlength{\itemsep}{0pt}\setlength{\parskip}{0pt}}
\setcounter{secnumdepth}{-\maxdimen} % remove section numbering
\usepackage{amsmath,amssymb,amsthm,mathrsfs}
\newcommand{\CPsh}{\operatorname{CPsh}}
\emergencystretch=3em
\ifLuaTeX
  \usepackage{selnolig}  % disable illegal ligatures
\fi
\usepackage{bookmark}
\IfFileExists{xurl.sty}{\usepackage{xurl}}{} % add URL line breaks if available
\urlstyle{same}
\hypersetup{
  pdftitle={Compact selection and a dual test-measure reduction for P21.3},
  pdfauthor={Xean native Codex research},
  hidelinks,
  pdfcreator={LaTeX via pandoc}}

\title{Compact selection and a dual test-measure reduction for P21.3}
\author{Xean native Codex research: gpt-6.1-sol, xhigh}
\date{3 October 2026}

\begin{document}
\maketitle

Independently reviewed supporting results, 3 October 2026. A separate
Xean native Codex review using gpt-6.1-sol at xhigh returned PASS for
the two support results, their first-stage corollary and fixed-current
estimate. The universal P21.3 remains unresolved: these are not a new
Bedford--Taylor positive range or an analytic counterexample. The full
review report and frozen proof accompany this note. No priority or formal
verification claim is made.

\section{1. Target, scope, and inputs}\label{target-scope-and-inputs}

The target fixes a complete non-Archimedean field \(K\), possibly
trivially valued, a good separated strict boundaryless pure
\(n\)-dimensional \(K\)-analytic space \(X\), \(n\geq2\), an arbitrary
\textbf{ordinary} open \(W\subset X\), and a positive, \(d'\)- and
\(d''\)-closed ordinary CLD current \(T\) of bidegree \((p,p)\),
\(0\leq p\leq n-1\). Its assertion is

\[
d'd''[U]=d'd''[V]\quad\Longrightarrow\quad
d'd''((U-V)T|_W)=0
\]

for every finite continuous CLD-psh pair \(U,V\) on \(W\). Here
\([f]=f\delta_X\), using the canonical integration current, including
generic component lengths. Positivity means the CLD middle cone.

This draft proves two support results.

\begin{enumerate}
\def\labelenumi{\arabic{enumi}.}
\tightlist
\item
  Transfinite decreasing constructions \textbf{on one unchanged domain},
  with no intermediate sheafification, have a compact upper-selection
  property. Their finite continuous outputs are locally decreasing
  uniform limits of smooth psh functions on a neighbourhood fixed before
  the sequence. Discontinuous intermediate functions cause no problem in
  this result.
\item
  For a fixed target current, annihilation of \textbf{all continuous
  zero-current functions} is equivalent to a precise weak closure
  condition on signed test measures. This is an exact dual reformulation
  of a stronger multiplier property, not a proof that the condition
  holds.
\end{enumerate}

Neither result assumes nontrivial valuation. Neither recognizes
arbitrary continuous positive-Hessian functions as CPsh, constructs
curvature-matched approximants for a pair, or supplies a density theorem
for currents.

The accepted strict real-PL theorem and its full proof and audit were
read in the supplied deliverables. Their scope is complete
\textbf{nontrivially} valued fields, actual GJR real-PL differences and
their stated fixed-neighbourhood locally uniform zero-current net
closure. That theorem is not a premise of either proof here. Its review
does not certify these new results or a changed recognition theorem. The
prior elliptic-disc surface, coarse fan, scalar-ascent, and
strong-product attempts were read and are not counted as new
counterexamples.

The primary source is {[}CLD{]}, arXiv:1204.6277v2, 25 July 2025. The
exact inputs needed here are:

\begin{itemize}
\tightlist
\item
  Definition 8.1.1: Psh is the smallest ordinary subsheaf of functions
  with values in \(\mathbb R\cup\{-\infty\}\) containing smooth
  functions with positive smooth Hessian and stable under decreasing
  sequential limits. The printed definition itself has no
  local-integrability clause; Proposition 8.1.12 separately treats locally
  integrable sections. The selection argument also works if local
  integrability is required at every limit, since it only uses an
  existing construction.
\item
  Definition 8.1.7: CPsh is the smallest ordinary subsheaf of finite
  continuous functions containing the same smooth generators and stable
  under uniform sequential limits. Lemmas 8.1.4 and 8.1.8 describe the
  respective \textbf{sheaf} stages: local decreasing/uniform limits at
  successors and local earlier-stage membership at limits.
\item
  Section 11.1.7: an analytic space has a basis of compact
  neighbourhoods; good spaces have affinoid neighbourhood bases. In the
  separated scope the underlying space is Hausdorff. Thus a point of an
  ordinary open has a compact neighbourhood contained in that open,
  whose ordinary interior still contains the point.
\item
  Theorem 15.7.1 and Section 15.7.2: a topologically separated pure
  analytic space has the canonical calibrated tropical structure of its
  dimension, and restriction to an analytic domain agrees with its own
  structure. Boundaryless ordinary opens are locally differentially
  separated. No valuation or reducedness restriction is imposed.
\item
  Proposition 6.5.2: a symmetric compactly supported smooth \((q,q)\)
  test is a difference of two middle-positive compactly supported smooth
  tests, provided its support has the stated paracompact and locally
  differentially separated neighbourhood. Its analytic support
  hypothesis holds for compact supports on an almost-good space, hence
  in the good boundaryless scope here.
\item
  Sections 7.2.8, 7.2.11--12: a positive current paired with a
  complementary positive smooth test gives a positive Radon coefficient
  measure; locally bounded Borel functions have the ordinary order-zero
  product. We give the needed signed-measure construction below using
  middle positivity and the preceding positive-test decomposition. We do
  not use a claim that strongly positive tests span all symmetric tests.
\item
  The ordinary current differential is dual to the test differential;
  for symmetric equal-bidegree tests the convention gives
  \(\langle d'd''(hT),\alpha\rangle
  =\langle hT,d'd''\alpha\rangle\). This is the convention used in Lemma
  7.7.5 and Section 8.5.3. The calculus has \(d'^2=d''^2=0\),
  \(d'd''=-d''d'\), and \(Jd'd''=d'd''J\).
\end{itemize}

The compact-selection proof below is elementary and proves its own Dini
statement. The dual-closure proof uses a finite-dimensional separation
argument given in full. Proposition 8.1.11's disputed reverse inclusion
and its same-rank assertion are not imported.

\section{2. Compact upper selection for upper semicontinuous
sequences}\label{compact-upper-selection-for-upper-semicontinuous-sequences}

\textbf{Lemma 1.} Let \(C\) be a compact topological space. Let
\(v_j:C\to\mathbb R\cup\{-\infty\}\) be upper semicontinuous and
decrease pointwise to \(v\). Let \(F:C\to\mathbb R\) be continuous and
assume \(v<F\) everywhere. Then some \(j\) satisfies \(v_j<F\)
everywhere on \(C\).

\textbf{Proof.} Each \(E_j=\{v_j<F\}\) is open in \(C\). They increase
with \(j\) and their union is \(C\), because of the strict pointwise
inequality and convergence. A finite subcover exists. Taking the largest
of its finitely many indices gives a single \(E_j=C\). This proof needs
neither metrizability nor continuity of the \(v_j\). \(\square\)

\textbf{Corollary 1 (upper semicontinuous Dini).} If the same sequence
has a finite continuous limit \(u\), it converges uniformly on \(C\).

\textbf{Proof.} Decrease gives \(v_j\geq u\). Apply Lemma 1 to
\(F=u+\varepsilon\) for each \(\varepsilon>0\). Eventually
\(0\leq v_j-u<\varepsilon\) on \(C\). This is precisely uniform
convergence. \(\square\)

Thus discontinuity of decreasing successor approximants does not itself
invalidate the Dini conclusion for a continuous finite output. However,
being a uniform approximant is not being a \textbf{CPsh} uniform
approximant. That distinction matters in the full sheaf construction.

\section{3. A construction with one fixed
domain}\label{a-construction-with-one-fixed-domain}

Fix a locally compact Hausdorff space \(O\) and a family
\(\mathcal S(O)\subset C(O,\mathbb R)\) of generators. Define classes of
functions with values in \(\mathbb R\cup\{-\infty\}\) by

\[
\begin{aligned}
\mathcal M_0(O)&=\mathcal S(O),\\
\mathcal M_{\alpha+1}(O)&=\mathcal M_\alpha(O)\cup
 \{\inf_j v_j: v_j\in\mathcal M_\alpha(O),\ v_{j+1}\leq v_j\},\\
\mathcal M_\lambda(O)&=\bigcup_{\beta<\lambda}\mathcal M_\beta(O)
 \quad (\lambda\text{ a nonzero limit ordinal}).
\end{aligned}
\]

Write \(\mathcal M(O)\) for their eventual union. This is a construction
in the set of functions on the fixed set \(O\); equivalently it is the
smallest class containing the generators and stable under decreasing
sequences. One may use the stages up to their stabilization: a strictly
increasing chain cannot add a new function at every stage beyond the
cardinality of this function set. In particular every member has an
ordinal rank. There is \textbf{no} gluing operation at successor stages
and \textbf{no} local earlier-rank membership at limit stages.

For the analytic generators, each \(\mathcal M_\alpha(O)\) is contained
in Psh\((O)\): the claim holds at rank zero, the Psh closure axiom
handles successors, and unions handle limit ranks. This inclusion does
not identify the fixed-domain class with the generated sheaf.

Every member is upper semicontinuous. This follows by transfinite
induction: the generators are continuous, a decreasing limit is upper
semicontinuous since \(\{\inf_jv_j<a\}=\bigcup_j\{v_j<a\}\), and limit
stages take unions of classes, not a new limit of functions.

\textbf{Theorem 1 (compact selection at arbitrary fixed-domain rank).}
For every \(v\in\mathcal M(O)\), every compact \(C\subset O\), and every
continuous \(F:C\to\mathbb R\) with \(v|_C<F\), there is a generator
\(s\in\mathcal S(O)\) such that

\[
v\leq s\text{ on all of }O,\qquad s|_C<F.
\]

\textbf{Proof.} Induct on the stage containing \(v\). At rank zero use
\(s=v\). At a successor, a member already in the previous class is
covered by induction. Otherwise it is the decreasing limit of
\(v_j\in\mathcal M_\alpha(O)\). Lemma 1 selects an index \(j\) with
\(v_j<F\) on the \textbf{same} compact \(C\). The induction hypothesis
selects \(s\in\mathcal S(O)\) satisfying \(v_j\leq s\) on all of \(O\)
and \(s<F\) on \(C\). Since \(v\leq v_j\) everywhere, this \(s\) works.
At a limit stage, \(v\) belongs on all of \(O\) to one earlier class;
apply that case. The induction never shrinks \(O\) or \(C\). This proves
the statement for every rank. \(\square\)

Notice that the limit at an intermediate node need not be continuous.
The continuous function is the strict upper obstacle \(F\) on the
compact. The selection of a single child does not assert uniform
convergence to a discontinuous intermediate limit.

\textbf{Theorem 2 (continuous output collapses to the first local
stage).} Let \(O\) be an ordinary open of the target analytic space.
Take \(\mathcal S(O)\) to be the ordinary CLD-smooth functions with
positive smooth Hessian on \(O\). Suppose \(u\in\mathcal M(O)\) is
finite and continuous. For every point \(x\in O\), choose
\textbf{before} the approximation index a compact neighbourhood
\(C\subset O\), and put \(B=\operatorname{Int}_O(C)\), with \(x\in B\).
There is a sequence of smooth psh functions on this one \(B\) which
decreases to \(u|_B\) and converges uniformly there. In particular
\(u\in\mathrm{CPsh}(O)\) and locally belongs to both first stages
Psh\(_1\) and CPsh\(_1\).

\textbf{Proof.} Set \(\varepsilon_m=3^{-m}\). Theorem 1 with \(v=u\) and
\(F=u|_C+\varepsilon_m\) gives a smooth psh \(s_m\) defined on
\textbf{all of the original} \(O\) with

\[
u\leq s_m\text{ on }O,\qquad
0\leq s_m-u<\varepsilon_m\text{ on }C.
\]

Define \(q_m=s_m+\varepsilon_m\), and restrict to \(B\). Addition of a
constant does not change its smooth Hessian. On \(B\),

\[
u+\varepsilon_m\leq q_m<u+2\varepsilon_m,
\qquad
q_{m+1}<u+2\varepsilon_{m+1}
 <u+\varepsilon_m\leq q_m.
\]

Thus \(q_m\) decreases on this fixed \(B\) and
\(\sup_B|q_m-u|\leq2\varepsilon_m\to0\). The first-stage conclusions
follow from the definitions; minimal CPsh and its ordinary sheaf
property give the result on \(O\). No positivity was obtained by a
partition of unity. No sequence of approximation-dependent domains or
intersection of infinitely many neighbourhoods occurs. \(\square\)

For precision one can sheafify \textbf{once, after} the fixed-domain
construction: let \(\mathcal M_{\mathrm{loc}}(W)\) consist of functions
which locally belong to \(\mathcal M(O)\) on an ordinary open
\(O\subset W\). Restriction sends each member of \(\mathcal M_\alpha(O)\)
into \(\mathcal M_\alpha(N)\) for every open \(N\subset O\),
by transfinite induction. Therefore this
local membership rule is a sheaf.

\textbf{Corollary 2.} As ordinary sheaves of finite continuous
functions,

\[
C\cap\mathcal M_{\mathrm{loc}}
 =\mathrm{CPsh}_1=C\cap\mathrm{Psh}_1.
\]

\textbf{Proof.} Theorem 2 gives the first inclusion into both first
stages. For the reverse inclusion from CPsh\(_1\), work on a
presentation open \(N\) where smooth psh functions converge uniformly to
\(u\). Choose members \(s_m\) with \(\sup_N|s_m-u|<\delta_m\),
\(\delta_m=4^{-m}\), and set \(r_m=s_m+2\delta_m\). Then

\[
u+\delta_m<r_m<u+3\delta_m,
\qquad r_{m+1}<u+3\delta_{m+1}<u+\delta_m<r_m.
\]

They decrease to \(u\) on \(N\). Thus \(u\in\mathcal M_1(N)\) and
\(u\in\mathrm{Psh}_1(N)\). Conversely a continuous Psh\(_1\)
presentation is a decreasing sequence of smooth functions on one open;
Corollary 1, or Theorem 2, gives a uniform sequence on a fixed
compact-neighbourhood interior. Ordinary locality yields the equality
claimed. \(\square\)

This does \textbf{not} identify \(\mathcal M_{\mathrm{loc}}\) with the
full Psh sheaf. In the full Psh construction a preceding-rank section
can have only local presentations, on different opens. A compact \(C\)
selected at the root may not lie in any one of them. Theorem 1 then
cannot descend to a single child defined on all of that \(C\). A finite
cover gives several local generators, not one positive global generator.
We have proved no positive gluing theorem for those generators. Nor have
we shown that the once-sheafified class is stable under all decreasing
sequences. No same-rank version of CLD 8.1.11 is asserted.

\section{4. Signed test measures and the stronger continuous kernel
property}\label{signed-test-measures-and-the-stronger-continuous-kernel-property}

We work now on an arbitrary ordinary \(W\) in the target, with the fixed
current \(T\). Let \(C(W)\) be all finite real continuous functions,
including those not globally bounded. Let \(\mathscr M_c(W)\) be the
real vector space of finite signed Radon measures with compact support.
Each such measure integrates every member of \(C(W)\). Equip this vector
space with the weak topology \(\sigma(\mathscr M_c(W),C(W))\): its
neighbourhoods are prescribed by finitely many continuous function
pairings. This is not convergence only against smooth functions.

If \(\beta\) is a symmetric compactly supported smooth
\((n-1,n-1)\)-form, define a signed measure \(\nu_\beta\) by

\[
\int_W f\,d\nu_\beta
 =\langle f\delta_X,d'd''\beta\rangle
 =\langle d'd''[f],\beta\rangle,\qquad f\in C(W).
\]

If \(\alpha\) is a symmetric compactly supported smooth
\((n-p-1,n-p-1)\)-form, similarly define \(\nu_{T,\alpha}\) by

\[
\int_W f\,d\nu_{T,\alpha}
 =\langle fT,d'd''\alpha\rangle
 =\langle d'd''(fT),\alpha\rangle.
\]

Here and below all current products are the ordinary products, and all
integration multiplicities belong to the original analytic calibration.

\textbf{Lemma 2 (existence of the signed test measures).} These formulas
define measures in \(\mathscr M_c(W)\), supported respectively in
\(\operatorname{supp}\beta\) and \(\operatorname{supp}\alpha\).

\textbf{Proof.} Let \(\eta\) stand for the complementary test
\(d'd''\alpha\); it is symmetric. Proposition 6.5.2 gives
\(\eta=\eta_1-\eta_2\), where \(\eta_i\) are positive compact tests.
Middle positivity of \(T\) makes \(g\mapsto\langle T,g\eta_i\rangle\) a
positive scalar distribution. A smooth cutoff equal to one on the
relevant compact gives its sup-norm bound. Smooth scalar density (CLD
6.2.5--6) extends it uniquely to a positive Radon measure, which is the
coefficient measure of CLD 7.2.11. Its mass is
\(\langle T,\eta_i\rangle\).

Their difference is a finite signed compact measure; pairing it with any
continuous \(f\) is the ordinary product of CLD 7.2.12. This definition
is independent of the decomposition: the pairings agree for all smooth
scalars, and smooth scalar density makes the Radon measures identical.
Outside \(\operatorname{supp}\eta\), smooth scalar pairings vanish; the
same density there proves that the signed measure has zero restriction.
Therefore its support is contained in
\(\operatorname{supp}\eta\subset\operatorname{supp}\alpha\). The case
\(T=\delta_X\) is the same argument. The source support hypotheses hold
on the good boundaryless ordinary open, as stated in Section 1. No
spanning assertion for strongly positive forms is used. \(\square\)

Write

\[
\mathscr L_W=\operatorname{span}_{\mathbb R}\{\nu_\beta:\beta
 \text{ as above}\},\qquad
\mathscr H_W=\{h\in C(W):d'd''[h]=0\}.
\]

By symmetry of the function Hessian and the current convention,
\(\mathscr H_W\) is exactly the annihilator of \(\mathscr L_W\) in
\(C(W)\). Testing symmetric forms suffices; antisymmetric forms pair to
zero.

In fact the ambient test space is linear and
\(\nu_{\beta_1}+a\nu_{\beta_2}=\nu_{\beta_1+a\beta_2}\), so its image
already equals the displayed real span. The span notation emphasizes
that no positivity restriction is imposed on these approximating
measures.

\textbf{Theorem 3 (dual reduction).} For the fixed \(T\), the following
are equivalent:

\begin{enumerate}
\def\labelenumi{\arabic{enumi}.}
\item
  For every \(h\in\mathscr H_W\), \(d'd''(hT)=0\).
\item
  For every test \(\alpha\) of the displayed degree,

  \[
  \nu_{T,\alpha}\in
  \overline{\mathscr L_W}^{\,\sigma(\mathscr M_c(W),C(W))}.
  \]
\end{enumerate}

\textbf{Proof.} If condition 2 holds, the continuous evaluation
\(\nu\mapsto\int h\,d\nu\) vanishes on \(\mathscr L_W\) for
\(h\in\mathscr H_W\), hence on its weak closure. It therefore vanishes
on every \(\nu_{T,\alpha}\), which is condition 1.

Conversely suppose a particular \(\nu=\nu_{T,\alpha}\) is outside that
closure. By the definition of the weak topology there are finitely many
\(f_1,\ldots,f_m\in C(W)\) and a positive tolerance for which the
corresponding basic neighbourhood of \(\nu\) misses \(\mathscr L_W\).
Define

\[
P:\mathscr M_c(W)\longrightarrow\mathbb R^m,
\qquad P(\mu)=(\int f_i\,d\mu)_{i=1}^m.
\]

Then \(P(\nu)\notin P(\mathscr L_W)\), for equality would produce a
measure of \(\mathscr L_W\) inside the forbidden neighbourhood. The
image \(P(\mathscr L_W)\) is a vector subspace of the finite-dimensional
space \(\mathbb R^m\), hence closed. There is a vector \(c\) orthogonal
to that subspace with \(c\cdot P(\nu)\ne0\): take the nonzero orthogonal
projection of \(P(\nu)\) onto the complement. Put \(h=\sum_i c_i f_i\).
It is finite continuous, annihilates every member of \(\mathscr L_W\),
and does not annihilate \(\nu\). Thus \(d'd''[h]=0\) but
\(\langle d'd''(hT),\alpha\rangle\ne0\), contradicting condition 1. This
proves the converse without importing a separation or current-density
theorem. \(\square\)

The closure may require a \textbf{net} of finite real linear
combinations of ambient Hessian test measures. It is not a
positive-current approximation and need not have common compact support
or a uniform total-variation bound. The convergence topology tests all
of \(C(W)\); convergence only against smooth scalars would not prove the
stated implication for a continuous \(h\). The theorem does not assert
that any of these nets exists for arbitrary \(T\).

Condition 1 is stronger than the target's property: the target uses only
\(h=U-V\) with finite continuous psh \(U,V\). Failure of condition 2
would supply a continuous analytic zero-current multiplier defect, but
the finite-dimensional separation proof does \textbf{not} make its \(h\)
a difference of finite continuous psh functions. Such failure alone
would therefore still not be an admissible P21.3 counterexample. No
failure of condition 2 is established in this draft.

\section{5. What is still missing for a
pair}\label{what-is-still-missing-for-a-pair}

The new selection argument concerns individual functions. Even if a full
Psh-to-CPsh conversion were later repaired, for \(h=U-V\) it would give
\(U,V\in\mathrm{CPsh}\), not \(h,-h\in\mathrm{CPsh}\). CPsh is a cone.
No operation in its definition subtracts two positive approximations
while preserving a curvature identity. The equality
\(d'd''[U]=d'd''[V]\) is upstairs against \(\delta_X\); it supplies no
equality of the finitely indexed approximants or of their products with
an arbitrary singular \(T\).

For a fixed \(T\) and fixed test \(\alpha\), Lemma 2 also gives the
legitimate estimate

\[
|\langle d'd''((f-g)T),\alpha\rangle|
 \leq \|f-g\|_{C^0(C)}\,
 \langle T,\eta_1+\eta_2\rangle,
\]

where \(C\) contains the positive-test decomposition supports for
\(d'd''\alpha=\eta_1-\eta_2\). The compact and coefficient are fixed
before the function approximation index. Consequently uniform scalar
approximation \textbf{does} control these differentiated pairings
against a fixed test. What it does not do is identify the limit as zero.
If \(h_j\to h\) locally uniformly, it shows
\(d'd''(h_jT)\to d'd''(hT)\). An additional argument is needed to show
that this limit vanishes. Distributional \(d'd''[h_j]\to0\) against the
ambient integration current supplies no comparison with the signed
measures \(\nu_{T,\alpha}\).

This distinction prevents a false diagnostic: when \(h_j\to0\) locally
uniformly, the displayed estimate already forces \(d'd''(h_jT)\to0\).
There cannot be a nonzero fixed-test defect in that situation. The
missing case has a potentially nonzero continuous zero-current limit
\(h\), not a zero scalar limit.

One possible independent route is now expressed exactly by Theorem 3:
prove its dual closure for the actual analytic test measures. No coarse
tropical pushforward, scalar extension of an arbitrary current, strong
current product, or weak-current density is a premise for that route.

\section{6. Self-audit and comparison with the residual
atom}\label{self-audit-and-comparison-with-the-residual-atom}

Theorem 1 uses a continuous strict upper obstacle on one fixed compact
and preserves the lower inequality on the entire domain. It does not use
false uniform convergence to a discontinuous intermediate limit. At
limit ranks its earlier-rank membership is global by definition;
substituting local membership would invalidate that proof step. Theorem
2 chooses its compact neighbourhood before \(m\), has smooth positive
Hessian at every finite step, and explicitly makes a decreasing
sequence. The one-time sheafification in Corollary 2 is distinguished
from sheafification between ordinal stages. These qualifications are
mathematical hypotheses, not cosmetic language.

Theorem 3 keeps the current arbitrary and treats every target degree,
including scalar tests at \(p=n-1\). Positive-test decomposition and
ordinary coefficient measures retain calibrations and generic lengths.
Its weak topology is spelled out, its necessity proof constructs a
continuous witness only, and no psh difference is inferred from that
witness. Neither theorem uses an unproved external recognition result.
They apply over trivial as well as nontrivial valuations, but provide no
new universal trivial-valuation Bedford--Taylor result.

The full Psh sheaf conversion still lacks a positive gluing argument on
a neighbourhood fixed before the tolerance. The universal kernel step
still lacks either two-sign recognition, a curvature-compatible
fixed-domain approximation/estimate, the dual closure of Theorem 3, or
another direct multiplier argument. No analytic counterexample
satisfying the equal-curvature and psh-pair premises was found. The
already accepted real-PL result is neither repeated as the new theorem
nor counted as a solution of the residual atom. The status is
\textbf{partial supporting progress; P21.3 unresolved}.

\section{7. Reproducibility and
bibliography}\label{reproducibility-and-bibliography}

All new outputs are in this worker directory. The standalone paper
\texttt{P21.3-partial.tex} contains the full proof bodies. Frozen inputs
and their SHA256 hashes are in \texttt{sources/} and
\texttt{source-manifest.json}; \texttt{sources/checked-passages.txt}
contains the actual CLD source extracts. \texttt{research-notes.md}
records the failed universal routes and concrete gaps.
\texttt{status.json} is the sole target status.
\texttt{environment.json} records the execution environment.
\texttt{self-audit.json} separates proof obligations from arithmetic
diagnostics.

Rerun \texttt{bash\ rerun.sh}. It checks frozen input hashes, runs the
small exact diagnostics in \texttt{diagnostics.py}, rebuilds TeX from
this argument using \texttt{build-paper.py}, and compiles twice using
the already installed TeX. The diagnostics and compiler output do not
establish mathematical correctness. No resource-heavy calculation,
installation, network call, subagent, or modification of earlier
campaigns was used. No operator-provided supervised compute service was
exposed.

{[}CLD{]} A. Chambert-Loir and A. Ducros, \emph{Formes différentielles
réelles et courants sur les espaces de Berkovich}, arXiv:1204.6277v2,
2025. \url{https://arxiv.org/abs/1204.6277v2}.

{[}GJR, contextual input only{]} W. Gubler, P. Jell and J. Rabinoff,
\emph{Forms on Berkovich spaces based on harmonic tropicalizations},
arXiv:2111.05741v3, 2025. \url{https://arxiv.org/abs/2111.05741v3}. The
accepted theorem's GJR real-PL scope is read as supplied context, not
used to prove these lemmas.

\end{document}
